% !PRINTED 564
\weberpage{564}
We shall not insert here the exact value for $K_n$, but rather a
smaller value that can be expressed more simply, whereby the
inequality (11) becomes all the more correct.

By the definition we have

\weq{17}{K_n = \int\!\int\!\cdots\!\int d\xi_1\, d\xi_2 \ldots d\xi_n,}

with the boundary

\weq{18}{\xi_1^2 + \xi_2^2 + \cdots + \xi_n^2 \leqq 1.}

The boundary equation (18) is certainly satisfied if we restrict
the $\xi_i$ to the intervals

\weq{19}{-\frac{1}{\sqrt{n}} < \xi_i < \frac{1}{\sqrt{n}}}

which geometrically has the meaning that we replace the sphere by
the inscribed cube. Accordingly we have

\[
K_n > \int\limits_{-\frac{1}{\sqrt n}}^{\frac{1}{\sqrt n}}
\int\limits_{-\frac{1}{\sqrt n}}^{\frac{1}{\sqrt n}}
\cdots
\int\limits_{-\frac{1}{\sqrt n}}^{\frac{1}{\sqrt n}}
d\xi_1\, d\xi_2 \cdots d\xi_n,
\]

and when this integration is carried out,

\weq{20}{K_n > \left(\frac{2}{\sqrt{n}}\right)^{n}.}

The upper limit determined in this way is never attained by $K_n$,
if we disregard the case $n = 1$.

The comparison of (20) with (16) yields a very simple expression
for the upper limiting value of $M$, namely

\weq{21}{M < n \sqrt[n]{D},}

and $M$ is (apart from the case $n = 1$) actually smaller than, and
not equal to, the value $n \sqrt[n]{D}$.

This shows itself clearly if one seeks a still somewhat narrower
limit for $M$.

The condition (18) is certainly then fulfilled if we set

\weq{22}{-\alpha < \xi_n < \alpha,\qquad \xi_1^2 + \xi_2^2 + \cdots + \xi_{n-1}^2 < 1 - \alpha^2}

where $\alpha$ may be an arbitrary positive proper fraction.
(The geometric meaning of this is that one replaces the sphere
by an inscribed cylinder of height $2\alpha$). Thus we have

\[
\begin{gathered}
K_n > 2\alpha \int\!\int\!\cdots\!\int d\xi_1\, d\xi_2 \ldots d\xi_{n-1}\\
\xi_1^2 + \xi_2^2 + \cdots + \xi_{n-1}^2 < 1 - \alpha^2.
\end{gathered}
\]
